\section{Proof of the exact normal-shell theorem}\label{ap:normal}

We prove Theorems~\ref{thm:normal-shell} and~\ref{thm:normal-taylor}.
Throughout, $m\ge3$, $n=m-1$, $\alpha=n/m$,
$J=\one\one^{\mathsf T}$, $P=I_m-J/m$, and $G_\Delta=P/\alpha$.
The reference configuration $V=[v_1\ \cdots\ v_m]$ satisfies
$V^{\mathsf T}V=G_\Delta$, and $d\ge m$.
Fix $\lambda>0$ and put $\mu=\lambda/\sqrt\alpha$.
The scalar results in Appendix~\ref{ap:transitions} give the strict normal
regime: $a_m(\mu)>m/2$ if $m\ge4$, and
$a_3(\mu)+b_3(\mu)>2$ if $m=3$.

We use coordinates on the positive square root of the centered Gram matrix.
They give a polynomial parametrization of the extended shell objective;
no differentiable choice of eigenvectors or implicit shell equation is needed.
The objective is extended to a full neighborhood in a fixed finite-dimensional
matrix space before it is restricted to realizable configurations.

\subsection{A smooth extension of the actual decoding deficit}\label{ng:smooth}

For a centered Gaussian vector $\xi$ with covariance $G$, the success
probability of a unit configuration is
\begin{equation}\label{ng:pc}
 \Pc(G;\lambda)=\frac{e^{-\lambda^2/2}}m
                 \E e^{\lambda\max_i\xi_i}.
\end{equation}
This formula is valid when $G$ is singular. If $Z$ is a standard Gaussian
in the signal space, the score vector
$(\langle x_i,Z\rangle)_{i=1}^m$ has exactly covariance $G$.

Let $b$ be an independent standard real Gaussian. The vector
$\sqrt\alpha\,\xi+m^{-1/2}b\one$ has covariance
$R=\alpha G+J/m$. Since $\mu\sqrt\alpha=\lambda$,
\[
 \E_R e^{\mu\max_i X_i}
   =e^{\mu^2/(2m)}\E_G e^{\lambda\max_i\xi_i},
 \qquad
 e^{-\mu^2/2}e^{\mu^2/(2m)}=e^{-\lambda^2/2}.
\]
Consequently the actual deficit is the restriction of
\begin{equation}\label{ng:extended-loss}
 \mathcal F_\lambda(R)=
 \frac{e^{-\mu^2/2}}m
 \left(\E_{I_m}e^{\mu\max_i Z_i}
             -\E_R e^{\mu\max_i X_i}\right):
 \qquad
 \mathcal F_\lambda(\alpha G+J/m)
   =\Pc(G_\Delta;\lambda)-\Pc(G;\lambda).
\end{equation}
The expectation defining this extension makes sense for every positive
definite symmetric $R$ near $I_m$, including matrices for which
$(R-J/m)/\alpha$ is not positive semidefinite.

Here is the regularity justification used below. On a closed covariance
neighborhood satisfying
\begin{equation}\label{ng:covariance-window}
 \tfrac12 I_m\preceq R\preceq2I_m,
\end{equation}
write
\[
 p_R(x)=(2\pi)^{-m/2}(\det R)^{-1/2}
                  \exp(-x^{\mathsf T}R^{-1}x/2).
\]
A derivative of order $k$ in the covariance entries is
$P_{k,R}(x)p_R(x)$, where $P_{k,R}$ is a polynomial of degree at most $2k$
whose coefficients are uniformly bounded on this neighborhood. Hence
\begin{equation}\label{ng:domination}
 |\partial_R^\beta p_R(x)|
 \le K_\beta(1+\norm{x}^{2|\beta|})e^{-\norm{x}^2/4}.
\end{equation}
Multiplication by $e^{\mu\max_i x_i}\le e^{\mu\norm{x}}$ preserves
integrability. Dominated differentiation therefore proves that
$\mathcal F_\lambda$ is smooth on a smaller neighborhood.
In particular it is $C^2$ jointly in all its covariance parameters, and
its third derivatives are locally bounded. All these neighborhoods and
bounds depend only on $m$ and $\lambda$.

For clarity, both bounds in \eqref{ng:covariance-window} are used:
the upper bound supplies the uniform Gaussian decay in
\eqref{ng:domination}. The argument estimates derivatives of the density,
rather than asserting a polynomial bound, uniform in $x$, on a density ratio.

\subsection{The first two covariance variations}\label{ng:variation}

Let $A=A^{\mathsf T}$ have zero diagonal, and set
\[
 S(A)=\sum_{i<j}A_{ij}=\tfrac12\one^{\mathsf T}A\one,
 \qquad N(A)=\norm{A}_F^2,\qquad R_1(A)=\norm{A\one}^2.
\]
We recall the scalar notation
\[
 q(c)=\frac{\phi(c)}{\Phi(c)},\qquad h(c)=c+q(c),\qquad
 \sigma(c)=1-q(c)h(c),\qquad k(c)=\phi(c)^2\Phi(c)^{m-2},
\]
and define the positive weight and its moments by
\begin{align}
 \omega_\lambda(c)&=\frac{\mu e^{-\mu^2/2}}m e^{\mu c}k(c),
       & C_\lambda&=\int_{\R}\omega_\lambda(c)\,dc,\label{ng:weight}\\
 A_\lambda&=\int_{\R}\omega_\lambda(c)h(c)^2\,dc,
       & B_\lambda&=\int_{\R}\omega_\lambda(c)\sigma(c)\,dc.\nonumber
\end{align}
Thus $C_\lambda>0$, $A_\lambda/C_\lambda=a_m(\mu)$, and
$B_\lambda/C_\lambda=b_m(\mu)$.

\begin{lemma}[Full covariance expansion]\label{ng:covariance-expansion}
For symmetric zero-diagonal $A$ near zero,
\begin{align}
 \mathcal F_\lambda(I_m+A)
 &=C_\lambda S(A)+\frac{A_\lambda}{4}N(A)
      -\frac{C_\lambda-B_\lambda}{2}R_1(A)\nonumber\\
 &\quad+\frac{S(A)^2}{2}
            \int_{\R}\omega_\lambda(c)q(c)^2\,dc
       +O(\norm{A}_F^3).
 \label{ng:full-expansion}
\end{align}
The remainder is uniform on a sufficiently small matrix neighborhood.
\end{lemma}

\begin{proof}
At the identity, the density Taylor coefficients are
\begin{align}
 p_{I_m+A}(x)
 =p_{I_m}(x)\bigl[
  1+\tfrac12x^{\mathsf T}Ax+\tfrac14N(A)
   +\tfrac18(x^{\mathsf T}Ax)^2
   -\tfrac12x^{\mathsf T}A^2x\bigr]+\mathcal R_A(x).
 \label{ng:density-expansion}
\end{align}
This follows by expanding $(I_m+A)^{-1}$ and
$-\tfrac12\log\det(I_m+A)$, using $\tr A=0$.
The density differentiation argument above gives, after shrinking the
neighborhood,
\[
 |\mathcal R_A(x)|
 \le K\norm{A}_F^3(1+\norm{x}^6)e^{-\norm{x}^2/4}.
\]
It thus also gives the integrated third-order remainder.

We identify the coefficients by integrating over $(-\infty,c]^m$.
The three one-dimensional truncated moments needed are
\[
 \int_{-\infty}^c\phi=\Phi(c),\qquad
 \int_{-\infty}^c x\phi(x)\,dx=-\phi(c),\qquad
 \int_{-\infty}^c x^2\phi(x)\,dx=\Phi(c)-c\phi(c).
\]
The linear coefficient of $A_{ij}$, $i<j$, is $k(c)$.
For the quadratic coefficients, an edge square $A_{ij}^2$ contributes
$c^2k(c)/2$; a product of distinct edges with one common vertex contributes
$-c q(c)k(c)$; and a product of disjoint edges contributes $q(c)^2k(c)$.
These assertions follow directly from \eqref{ng:density-expansion}:
the corresponding polynomials are
\[
 \tfrac12(1-x_i^2)(1-x_j^2),\qquad
 (x_i^2-1)x_jx_\ell,\qquad x_ix_jx_\ell x_s,
\]
respectively. Each product of distinct edges is counted once.

Put $Q_e=\sum_{i<j}A_{ij}^2=N(A)/2$, and denote the sums of adjacent
and disjoint edge products by $U_e$ and $W_e$. Then
\[
 R_1(A)=2Q_e+2U_e,\qquad
 S(A)^2=Q_e+2U_e+2W_e.
\]
Substitution gives the fixed-$c$ expansion
\begin{align}
 \frac{F_{I_m+A}(c)-\Phi(c)^m}{k(c)}
 =S(A)+\frac{h(c)^2}{4}N(A)
    -\frac{q(c)h(c)}2R_1(A)
    +\frac{q(c)^2}2S(A)^2+O_c(\norm{A}_F^3).
 \label{ng:orthant-expansion}
\end{align}

We justify the subsequent integration, including its sign. Expressing
the expectations through the distribution of the maximum gives
\begin{equation}\label{ng:layer-cake}
 \E_{I_m}e^{\mu\max_i Z_i}-\E_R e^{\mu\max_i X_i}
   =\mu\int_{\R}e^{\mu c}\bigl(F_R(c)-\Phi(c)^m\bigr)\,dc.
\end{equation}
One may obtain it by integration by parts, or by applying Fubini to the
difference of the two laws, which has total mass zero. The same argument
applies to a positive-order covariance derivative of the density:
its integral over all of $\R^m$ is zero. For $c\ge0$, its integral over
$(-\infty,c]^m$ is minus its integral over the complementary set, which is
contained in $\{\norm{x}\ge c\}$. For $c\le0$, the orthant itself lies in
$\{\norm{x}\ge|c|\}$. Formula~\eqref{ng:domination} therefore bounds each
first or second orthant derivative by
$K(1+|c|^N)e^{-c^2/8}$. This remains integrable after multiplication by
$e^{\mu c}$. It validates Fubini and differentiation of
\eqref{ng:layer-cake} through order two.

Multiply the explicitly displayed normalized coefficients in
\eqref{ng:orthant-expansion} by the full weight $\omega_\lambda(c)$,
which includes $k(c)$, and integrate. Since $qh=1-\sigma$, this yields
\eqref{ng:full-expansion}. The bound
$q(c)\le |c|+\sqrt{2/\pi}$ ensures finiteness of the $q^2$ coefficient.
The integrated remainder comes from the smooth density argument, not
from integrating an unspecified fixed-$c$ remainder.
\end{proof}

\subsection{Compression and its exact inverse}\label{ng:compression}

Identify operators on $H=\one^\perp$ with symmetric $m\times m$ matrices
$C$ satisfying $PC=CP=C$. Define the linear map
\begin{equation}\label{ng:uncompress}
 \mathcal U(C)_{ij}=C_{ij}-\tfrac12(C_{ii}+C_{jj}).
\end{equation}
It has zero diagonal, and direct multiplication gives
\begin{equation}\label{ng:inverse}
 P\mathcal U(C)P=C,\qquad
 \mathcal U(PAP)=A
 \quad\text{if }A=A^{\mathsf T},\ \diag A=0.
\end{equation}
For the second identity, write $a=A\one$ and
$s=\one^{\mathsf T}A\one$. The entries of $PAP$ are
$A_{ij}-(a_i+a_j)/m+s/m^2$, and its diagonal entries are
$-2a_i/m+s/m^2$; \eqref{ng:uncompress} cancels the extra terms.
Since a centered matrix has zero total entry sum,
\begin{equation}\label{ng:trace-sum}
 S(\mathcal U(C))=-\frac m2\tr C.
\end{equation}

The shape space used in this appendix is
\[
 \mathcal D_m=\{D=D^{\mathsf T}:PD=DP=D,\ \tr D=0\}.
\]
If $D\in\mathcal D_m$ and $B=\mathcal U(D)$, then $S(B)=0$ and $D=PBP$.
Writing $b=B\one$, we have $\one^{\mathsf T}b=0$ and
\[
 D=B-\frac1m(b\one^{\mathsf T}+\one b^{\mathsf T}).
\]
Expansion of the squared Frobenius norm gives the useful exact identity
\begin{equation}\label{ng:metric-identity}
 \tr(D^2)=\norm{D}_F^2
       =N(B)-\frac2mR_1(B).
\end{equation}

\subsection{The actual Procrustes infimum and a uniform neighborhood}
\label{ng:procrustes}

For unit configurations define
\[
 d_P(X,V)^2=\inf_{U\in O(d)}\norm{X-UV}_F^2.
\]
This is the actual configuration distance. The following calculation
both evaluates the infimum near a simplex and exhibits an attaining
orthogonal map.

Let $G=X^{\mathsf T}X$, and set
\begin{equation}\label{ng:root}
 T=(\alpha PGP)^{1/2},\qquad K=T+(I_m-P).
\end{equation}
The square root is positive semidefinite, $PT=TP=T$, and $T^2=\alpha PGP$.
Assume for the moment that $K$ is invertible. Put
\[
 W=XPK^{-1}.
\]
The regular Gram identity implies $V\one=0$.
The matrix $K$ preserves $H$ and its orthogonal complement; on these
spaces it acts as $T$ and the identity, respectively. Consequently
\[
 W^{\mathsf T}W
 =K^{-1}PGPK^{-1}=\alpha^{-1}P=G_\Delta,
 \qquad WP=W,\qquad XP=WT.
\]
Configurations with the same Gram matrix are related by an orthogonal
map of the ambient space. Indeed, sending the columns of one to those
of the other defines an isometry between their spans: the Gram equality
makes the map well-defined and preserves every inner product.
An orthogonal extension over the complementary subspaces gives an element
of $O(d)$. Thus $W=U_*V$ for some $U_*\in O(d)$.

For any other $Y=UV$, the matrices $Y$ and $W$ are centered and have Gram
matrix $P/\alpha$. Positivity of $T$ gives
\[
 0\le\tr\!\left(T(W-Y)^{\mathsf T}(W-Y)\right)
   =\frac2\alpha\tr T-2\tr(TW^{\mathsf T}Y).
\]
Because $XP=WT$ and $YP=Y$, the last trace equals
$\langle X,Y\rangle_F$. It is at most $\tr T/\alpha$,
with equality at $Y=W$. Since $\norm X_F^2=\norm V_F^2=m$, we obtain
\begin{equation}\label{ng:distance}
 d_P(X,V)^2=2m-\frac2\alpha\tr T.
\end{equation}
In particular the infimum is attained in the situation used here.

We now verify uniformly when this situation applies. If
$\norm{X-UV}_F^2=e$, then, for each pair of indices,
\[
 |G_{ij}-(G_\Delta)_{ij}|
 \le\norm{x_i-Uv_i}+\norm{x_j-Uv_j}\le2\sqrt e.
\]
Taking a sequence approaching the infimum gives
\begin{equation}\label{ng:gram-enclosure}
 \norm{G-G_\Delta}_{\mathrm{op}}
 \le\norm{G-G_\Delta}_F\le2m\,d_P(X,V).
\end{equation}
The bound is independent of $d$ and $V$.

On the cone of positive semidefinite $m\times m$ matrices, the maps
\[
 G\longmapsto(\alpha PGP)^{1/2}+(I_m-P),\qquad
 G\longmapsto\alpha G+J/m
\]
are continuous and take the value $I_m$ at $G_\Delta$.
Invertibility is an open condition. The second map is positive
semidefinite on this cone, so invertibility makes it positive definite.
There is therefore an $\varepsilon_m>0$, in this fixed matrix space,
such that both required conclusions hold when
$\norm{G-G_\Delta}_{\mathrm{op}}<\varepsilon_m$.
Together with \eqref{ng:gram-enclosure}, this supplies a positive
Procrustes neighborhood depending only on $m$. We have used continuity
of the square root on the positive cone, including at the singular
matrix $P$, but have not differentiated that square root.

\subsection{Polynomial square-root coordinates and enclosure}
\label{ng:chart}

Put $\tau=\rho^2$ and
\[
 \kappa=\kappa(\tau)=1-\frac{\tau}{2m}.
\]
If $X$ is on a sufficiently small actual shell $d_P(X,V)^2=\tau$,
\eqref{ng:distance} gives $\tr T=(m-1)\kappa$. Hence
\begin{equation}\label{ng:actual-deviation}
 D=T-\kappa P\in\mathcal D_m.
\end{equation}
Conversely, on the fixed vector space $\mathcal D_m$ define the polynomial
maps
\begin{equation}\label{ng:polynomial-chart}
 C_\tau(D)=(\kappa P+D)^2-P,\qquad
 R_\tau(D)=I_m+\mathcal U(C_\tau(D)),\qquad
 f_\tau(D)=\mathcal F_\lambda(R_\tau(D)).
\end{equation}
For $(\tau,D)$ near $(0,0)$, $R_\tau(D)$ is positive definite and
$f_\tau(D)$ is jointly $C^2$. This definition permits all small
$D\in\mathcal D_m$, irrespective of whether the corresponding Gram
matrix is realizable.

Every actual sufficiently small shell is represented by this chart.
Indeed, for its $D$ in \eqref{ng:actual-deviation},
\[
 C_\tau(D)=T^2-P=P(\alpha G-P)P.
\]
The matrix $\alpha G-P$ has zero diagonal. The inverse identity
\eqref{ng:inverse} gives the exact reconstruction
\begin{equation}\label{ng:reconstruction}
 R_\tau(D)=\alpha G+J/m.
\end{equation}
There is also a quantitative enclosure that uses no modulus of continuity
for the square root. Since the columns of $X$ are unit vectors,
\begin{align*}
 \tr(T^2)
 &=\alpha\tr(PGP)
   =\alpha\left(m-\frac{\norm{\sum_i x_i}^2}{m}\right)
   \le m-1,\\
 \norm D_F^2
 &=\tr(T^2)-\kappa^2(m-1)
   \le(m-1)(1-\kappa^2)\le\tau.
\end{align*}
Thus
\begin{equation}\label{ng:whole-shell-enclosure}
 \norm D_{\mathrm{op}}\le\norm D_F\le\rho
\end{equation}
for every configuration on every sufficiently small positive shell.
The entire feasible shell, not just a chosen smooth path through the
simplex, is covered.

\subsection{Stationarity and the exact shell Hessian}\label{ng:hessian}

Conjugation by a permutation matrix preserves $P$, $\mathcal D_m$,
the map $\mathcal U$, and the chart \eqref{ng:polynomial-chart}.
If $X\sim N(0,R)$ and $Q$ is a permutation matrix, then
$QX\sim N(0,QRQ^{\mathsf T})$, while $\max_i(QX)_i=\max_iX_i$.
Together with the equivariance of the chart, this proves that $f_\tau$
is permutation invariant.
For each sufficiently small fixed $\tau$, its differential at zero is
an invariant linear functional on $\mathcal D_m$. The average of any
$D\in\mathcal D_m$ under all permutations is a matrix $aI_m+bJ$.
Centeredness gives $a+mb=0$, while zero trace gives $a+b=0$.
Since $m>1$, both coefficients vanish. Applying the invariant functional
to this average proves
\begin{equation}\label{ng:stationarity}
 \mathrm d f_\tau(0)[D]=0\qquad(D\in\mathcal D_m).
\end{equation}

We compute the Hessian at $(\tau,D)=(0,0)$.
Fix $D\in\mathcal D_m$ and write $B=\mathcal U(D)$ and
$L=\mathcal U(D^2)$. The polynomial chart gives exactly
\[
 R_0(tD)=I_m+2tB+t^2L.
\]
Here $S(B)=0$ and, by \eqref{ng:trace-sum},
$S(L)=-m\tr(D^2)/2$.
Substitution in \eqref{ng:full-expansion} yields
\begin{equation}\label{ng:Q}
 f_0(tD)=t^2 Q(D)+O(t^3),\qquad
 Q(D)=-\frac{mC_\lambda}{2}\tr(D^2)
      +A_\lambda N(B)-2(C_\lambda-B_\lambda)R_1(B).
\end{equation}
The $S(A)^2$ term of the full expansion is of order $t^4$ here.
It is needed in the general covariance expansion but contributes
nothing to this Hessian. Since $f_0$ is $C^2$,
\begin{equation}\label{ng:actual-hessian}
 \mathrm d^2 f_0(0)[D,D]=2Q(D).
\end{equation}
This follows either by the second derivative chain rule or by comparing
the $t^2$ Taylor coefficient. No remainder symbol is differentiated.

For comparison with the in-plane sectors, write $b=B\one$ and set
\[
 (B_1)_{ii}=0,\qquad
 (B_1)_{ij}=\frac{b_i+b_j}{m-2}\ (i\ne j),\qquad B_0=B-B_1.
\]
Since $\one^{\mathsf T}b=0$, these matrices satisfy
$B_0\one=0$, $B_1\one=b$, and $\langle B_0,B_1\rangle_F=0$.
Direct summation also gives
$\norm{B_1}_F^2=2\norm b^2/(m-2)$.
Define
\[
 E_0=\frac{\norm{B_0}_F^2}{4\alpha},\qquad
 E_1=\frac{\norm b^2}{m-2}.
\]
They are nonnegative and vanish together only when $D=0$. Moreover,
\[
 N(B)=4\alpha E_0+2E_1,\qquad R_1(B)=(m-2)E_1.
\]
For $m=3$, the symmetric zero-diagonal row-zero matrix $B_0$ is zero:
its three row equations force all three off-diagonal entries to vanish.
Thus $E_0=0$ in that case.

With
\[
 \Lambda_0=\alpha A_\lambda,\qquad
 \Lambda_1=\frac{A_\lambda+(m-2)B_\lambda}{2},\qquad
 \Lambda_\perp=\frac{m-1}{2}C_\lambda,
\]
identities \eqref{ng:metric-identity} and \eqref{ng:Q} give
\begin{equation}\label{ng:hessian-gaps}
 Q(D)=4\bigl[(\Lambda_0-\Lambda_\perp)E_0
               +(\Lambda_1-\Lambda_\perp)E_1\bigr],
 \qquad
 \mathrm d^2 f_0(0)[D,D]
   =8\bigl[(\Lambda_0-\Lambda_\perp)E_0
               +(\Lambda_1-\Lambda_\perp)E_1\bigr].
\end{equation}
The factors in this display are part of the coordinate normalization:
the covariance velocity in the present chart is $2B$, and the Hessian
is twice the quadratic Taylor coefficient.

The scalar gap assumptions make \eqref{ng:hessian-gaps} positive definite.
One can see the required quantitative coercivity without any additional
representation-theoretic argument. Formula~\eqref{ng:metric-identity}
reduces \eqref{ng:Q} to
\begin{equation}\label{ng:raw-coercivity}
 Q(D)=\left(A_\lambda-\frac m2C_\lambda\right)N(B)
                       +(2B_\lambda-C_\lambda)R_1(B).
\end{equation}
For $m\ge4$, Appendix~\ref{ap:transitions} gives
$b_m(\mu)>1/2$ for $\mu\ge0$, while the strict normal regime gives
$a_m(\mu)>m/2$. Both coefficients in \eqref{ng:raw-coercivity}
are positive, and $N(B)\ge\norm D_F^2$.
For $m=3$, $R_1(B)=N(B)/2$, so
\[
 Q(D)=(A_\lambda+B_\lambda-2C_\lambda)N(B),
\]
again with a positive coefficient. In either case there is a
$\gamma=\gamma(m,\lambda)>0$ such that
$Q(D)\ge\gamma\norm D_F^2$ for every $D\in\mathcal D_m$.

\subsection{Uniform convexity and the equality orbit}\label{ng:minimum}

Continuity of the Hessian of the joint $C^2$ map $(\tau,D)\mapsto f_\tau(D)$,
together with \eqref{ng:actual-hessian}, gives positive numbers
$\tau_*$, $\delta$, and $\eta$, depending only on $m$ and $\lambda$, such that
\[
 \mathrm d^2f_\tau(D)[H,H]\ge\eta\norm H_F^2
 \quad\text{if }|\tau|<\tau_*,\ \norm D_F<\delta,\
                       H\in\mathcal D_m.
\]
Shrink the neighborhoods if necessary to keep the covariance positive
definite. The ball in $\mathcal D_m$ is convex. Integrating the second
derivative along $t\mapsto tD$, and using \eqref{ng:stationarity}, gives
\begin{equation}\label{ng:strict-minimum}
 f_\tau(D)-f_\tau(0)
 =\int_0^1(1-t)\,\mathrm d^2 f_\tau(tD)[D,D]\,dt
 \ge\frac\eta2\norm D_F^2.
\end{equation}
In particular, $f_\tau(D)=f_\tau(0)$ holds if and only if $D=0$.

Choose $\rho_0>0$ smaller than $\sqrt{2m}$, $\sqrt{\tau_*}$,
$\delta$, and the uniform Procrustes neighborhood from
Subsection~\ref{ng:procrustes}. These choices involve only $m$ and
$\lambda$, before $d$ or $V$ is selected.
By \eqref{ng:whole-shell-enclosure}, every actual unit configuration on
a shell $0<\rho<\rho_0$ satisfies the hypotheses of
\eqref{ng:strict-minimum}. By \eqref{ng:reconstruction} and
\eqref{ng:extended-loss}, its value $f_{\rho^2}(D)$ is its actual
decoding deficit.

It remains to verify attainment and identify every equality case.
Since $\mathcal U(P)=I_m-J$, the point $D=0$ has
\[
 R_{\rho^2}(0)=\kappa^2I_m+(1-\kappa^2)J,\qquad
 G_\kappa=\kappa^2G_\Delta+(1-\kappa^2)J,
 \quad \kappa=1-\frac{\rho^2}{2m}\in(0,1).
\]
The span of $V$ has dimension $m-1$, so $d\ge m$ provides a unit
$z\perp\operatorname{span}\{v_i\}$. Define
\begin{equation}\label{ng:normal-realization}
 N_i=\kappa v_i+\sqrt{1-\kappa^2}\,z.
\end{equation}
These are unit vectors and have Gram matrix $G_\kappa$.
Their centered square root is $T=\kappa P$, so $K$ is invertible and
\eqref{ng:distance} gives
\[
 d_P(N,V)^2=2m(1-\kappa)=\rho^2.
\]
Thus the minimizer of the extended problem is feasible in the original
problem. Its Gram matrix has eigenvalues $\kappa^2/\alpha$ on $H$
and $m(1-\kappa^2)$ on $\operatorname{span}\{\one\}$, so it has rank $m$.

For an arbitrary configuration on this shell, equality in the deficit
bound forces $D=0$ by \eqref{ng:strict-minimum}, hence $G=G_\kappa$ by
reconstruction. The same-Gram isometry argument in
Subsection~\ref{ng:procrustes} then gives precisely the $O(d)$ orbit
of \eqref{ng:normal-realization}. Conversely every configuration on
this orbit has the same Gram matrix and success probability.
Different choices of the unit normal, and permutations of the regular
simplex vertices, produce the same orbit. This proves the equality
assertion and the uniform radius statement of
Theorem~\ref{thm:normal-shell}.

\subsection{The exact value and the sixth-order remainder}\label{ng:taylor}

Let $\xi\sim N(0,G_\Delta)$ and let $b$ be an independent standard real
Gaussian. The vector
$\kappa\xi+\sqrt{1-\kappa^2}\,b\one$ has covariance $G_\kappa$.
For $\kappa\ge0$, its coordinate maximum is
$\kappa\max_i\xi_i+\sqrt{1-\kappa^2}\,b$.
Using the Gaussian moment-generating function in \eqref{ng:pc} yields
\begin{equation}\label{ng:exact-value}
 \Pc(G_\kappa;\lambda)
   =\frac{e^{-\lambda^2/2}}m
      e^{\lambda^2(1-\kappa^2)/2}
      \E e^{\lambda\kappa\max_i\xi_i}
   =\Pc(G_\Delta;\lambda\kappa).
\end{equation}
This proves the exact optimal value, with all prefactors included.

Write $p(s)=\Pc(G_\Delta;s)$ using the expectation formula as a real
extension in $s$. This extension is smooth. Indeed, if
$M=\max_i\xi_i$, then each coordinate $\xi_i$ is Gaussian and
$|M|\le\max_i|\xi_i|$. Therefore
$\E(|M|^k e^{L|M|})<\infty$ for every fixed $k$ and $L$.
On compact $s$-intervals these moments dominate all the derivatives of
$e^{-s^2/2}\E e^{sM}/m$. In particular $p$ is $C^3$ near $\lambda$,
with bounds depending only on $m$ and $\lambda$.

Choose $u>0$ and $M_3>0$ such that
$|p'''(s)|\le M_3$ whenever $|s-\lambda|\le u$.
For $h=\lambda\rho^2/(2m)$ small enough, Taylor's theorem gives
\[
 \left|p(\lambda)-p(\lambda-h)
       -p'(\lambda)h+\frac12p''(\lambda)h^2\right|
       \le\frac{M_3}{6}|h|^3.
\]
Together with \eqref{ng:exact-value}, this proves
\begin{equation}\label{ng:O6}
 \delta_{\min}(\rho)
 =\frac{\lambda p'(\lambda)}{2m}\rho^2
    -\frac{\lambda^2p''(\lambda)}{8m^2}\rho^4+O(\rho^6).
\end{equation}
More explicitly, after decreasing $\rho_0$ if necessary there is a
$C>0$, depending only on $m$ and $\lambda$, such that
\begin{equation}\label{ng:quantified-remainder}
 \left|\delta_{\min}(\rho)
       -\frac{\lambda p'(\lambda)}{2m}\rho^2
       +\frac{\lambda^2p''(\lambda)}{8m^2}\rho^4\right|
 \le C|\rho|^6,\qquad 0<\rho<\rho_0.
\end{equation}
For example the Taylor argument permits $C=M_3\lambda^3/(48m^3)$
on a sufficiently small interval. The same estimate holds on a
two-sided interval for the real extension
$p(\lambda)-p(\lambda(1-\rho^2/(2m)))$.

That exact extension is even in $\rho$, so every odd derivative at zero
vanishes. In particular there is no cubic term; this follows from the
exact value, not from dropping an uncomputed cubic remainder.
The kernel-checked Taylor argument proves the same quantified
$O(\rho^6)$ statement by applying the order-three Taylor little-$o$
theorem, adding back its cubic coefficient to obtain
$O(h^3)$ after the quadratic polynomial, and composing with
$h=-\lambda\rho^2/(2m)$. It thus verifies the full sixth-order bound,
in addition to the parity identity.

The proof uses strict positivity of the existing in-plane curvature
gaps. At a threshold where one such gap vanishes,
\eqref{ng:strict-minimum} is not supplied by this argument.
No assertion about that degenerate case, or an effective numerical
lower bound for $\rho_0$, is part of the theorem.
